% !TEX root = ../main.tex
\lecture{4}{2026-10-05}{Doppelte Buchhaltung}

\section{Fazit: Konzept der doppelten Buchhaltung}

\subsection{Kontensystem}

\subsection{Kontenführung}

\begin{enumerate}
    \item Festlegung der relevanten Konten
    \begin{itemize}
        \item anspruchsvoll zu wissen welche relevant sind
        \item hilfreich: Kontenplan
        \item resultat: Hauptbuch
    \end{itemize}
    \item Übertrag der Eröffnungsbilanz in Aktiven und Passiven
    \item Erfassung aller relevanten Buchungstatsachen.
    \item Kontenführung
    \begin{itemize}
        \item Eintrag von Buchungssätze in Konten (1x Soll, 1x Haben)
    \end{itemize} 
    \item Abschluss 
    \begin{itemize}
        \item Zu Rechnung an einem Stichtag gemacht (z.B. 31.12)
        \item alle Konten Saldiert (SB bei Bilanz, und Saldi Erfolgsrechnung)
        \item Summe der SBs = Schlussbilanz
        \item ER: Summe der Saldi = Erfolgsrechnung
        \item In der ER (und in der Bilanz): Differenz von Gewinn und Verlust, diese Differenzen gleichen sich aus im Vergleich von Bilanz und ER
        \item Velust oder Gewinn geht ins \textbf{Eigenkapitalkonto}
        \item Beispiel
        \begin{itemize}
            \item Gewinn 50k
            \item 25k im unternehmen (Eigenkapital)
            \item 25k Ausschüttung (Dividenden, \textbf{Aktienrückkauf}) \footnote{Wie siehts in der Schweiz aus, Dividendesteuer etc?} \footnote{Das Kapital fliesst in der Richtung der grössten Profitraten. Woher und wie es überhaupt möglich ist aus dem Tausch, oder Kauf und Verkauf von Äquivalente Mehrwert zu schaffen, bleibt unberuht.}
        \end{itemize}
    \end{itemize}
    \item Analyse (Kapitel Kennzahlen)
\end{enumerate}

\pagebreak
\subsection{Abschluss Tabellenform}

\begin{enumerate}
    \item Probebilanz
    \begin{itemize}
        \item zeigt alle Beträge ohne Saldierung
        \item 
    \end{itemize}
    \item Saldobilanz
    \begin{itemize}
        \item zeigt nur SBs und Saldi
        \item zb Monatlich
    \end{itemize}
    \item Bilanz vor Gewinnverteilung
    \item Erfolgsrechnung
    \begin{itemize}
        \item Hier wird der Gewinn explizit angegeben, um die Differenz zur Bilanz v.Gv. auszugleichen.
    \end{itemize} 
    \item Bilanz nach Gewinnverteilung
    \begin{itemize}
        \item Gewinn wird über die Konten verteilt
        \item Meistens ins Eigenkapital oder Dividenden/Aktienrückkauf
        \item Buchungsatz bei Ausschüttung
        \subitem \textbf{Fremdkapital} / Kasse 200 (Definanzierung)
    \end{itemize} 
\end{enumerate}

\subsubsection{Buchungssätze: Kapitalquellen}

\begin{itemize}
    \item unternehmen
    \begin{itemize}
        \item Gewinn
        \begin{itemize}
            \item Aufwand / Kasse 
            \item Kasse / Ertrag 
        \end{itemize}
        \item Verkauf Aktiven 
        \begin{itemize}
            \item Kasse / Gebäude 
        \end{itemize}
    \end{itemize}
    \item Kapitalaufnahme
    \begin{itemize}
        \item Erhöhung Eigenkapital 
        \begin{itemize}
            \item Kasse / EK
        \end{itemize}
        \item Erhöhung Fremdkapital
        \begin{itemize}
            \item Kasse / Fremdkapital
        \end{itemize}
    \end{itemize}
\end{itemize}

\subsubsection{Buchungssätze: Kapitalgeber}

\begin{itemize}
    \item unternehmen
    \begin{itemize}
        \item Unternehmensaktisition 
    \end{itemize}
    \item Rückzahlung an Kapitalgeber
    \begin{itemize}
        \item Rückkäufe
        \begin{itemize}
            \item EK / Kasse
        \end{itemize}
        \item Dividendenzahlungen
        \begin{itemize}
            \item FK / Kasse
        \end{itemize}
        \item Rückzahlung FK
        \begin{itemize}
            \item FK /Kasse
        \end{itemize}
    \end{itemize}
\end{itemize}